
\def\PUDcodigo{ECA222}

\if\COMPILAR0
   \def\PUDcodacad{04507.30}
\else
   \def\PUDcodacad{04506.31}
\fi

\def\PUDdisciplina{Instalações Elétricas}

\def\PUDsemestre{S6}

\def\PUDhoras{80}

%NOVOS ---------------------------------------------------
\def\PUDhorasTeoria{60}
\def\PUDhorasPratica{20}

\if\COMPILAR0
   \def\PUDinterECA{ 
   \hyperref[PUD:EME106]{Desenho Auxiliado por Computador (04507.11)}
   
   \hyperref[PUD:ECA215]{Circuitos Elétricos II (04507.24)}
   
   \hyperref[PUD:ECA219]{Instrumentação (04507.36)} 
   }
\else
   \def\PUDinterECA{ 
   \hyperref[PUD:EME106]{Desenho Auxiliado por Computador (04506.10)}
   
   \hyperref[PUD:ECA210]{Circuitos Elétricos I (04506.21)}
   }
\fi

\def\PUDatuaisECA{
- Automação residencial e predial

- Automação de subestações

- Internet das coisas (IoT) aplicada a equipamentos residenciais, prediais e industriais
}

\def\PUDrecursos{Projetor multimídia; Quadro branco e pincel; Aulas em laboratório de eletricidade CA e visita às instalações do campus.}

%----------------------------------------------------------

\def\PUDcreditos{4}

\def\PUDprerequisito{ \hyperref[PUD:ECA210]{Circuitos Elétricos I (ECA210)} }

\if\COMPILAR0
   \def\PUDpreacad{ \hyperref[PUD:ECA210]{Circuitos Elétricos I (04507.20)} }
\else
   \def\PUDpreacad{ \hyperref[PUD:ECA210]{Circuitos Elétricos I (04506.21)} }
\fi

\def\PUDobjetivos{Conhecer as principais normas técnicas, como a NR10, a NBR 5410, a NBR 5419 e outras.  Identificar os equipamentos que fazem parte de uma instalação elétrica industrial.  Conhecer e especificar os dispositivos que compõem um sistema de iluminação industrial.  Especificar condutores elétricos para circuitos de iluminação e força industriais em baixa tensão.  Conhecer e especificar condutos.  Conhecer e especificar bancos de capacitores.  Conhecer e especificar os dispositivos fundamentais de proteção elétrica em baixa tensão.  Conhecer os sistemas de aterramento e para-raios.  Conhecer as subestações e geração elétrica industrial.}

\def\PUDementa{Fornecimento de Energia Elétrica à Indústria;  Tarifação da Energia Elétrica;  Fator de Potência e Banco de Capacitores;  Iluminação Industrial;  Dimensionamento de Condutores e Condutos;  Alimentação de Circuitos de Motores Elétricos;  Proteção e Coordenação;  Aterramento;  Proteção Contra Descargas Atmosféricas.}

\def\PUDconteudo{ & Sistemas Elétricos de Potência.  Níveis de Tensão em uma Instalação Elétrica Industrial.  Subestações Industriais.  Produção de Energia Elétrica na Indústria.  Cogeração. 
\\ 
 & Definições e Conceitos.  Níveis de Tensão e Estruturas Tarifárias.  Faturamento. 
\\ 
 & Potências em CA e Fator de Potência.  Correção do Fator de Potência.  Bancos de Capacitores. 
\\ 
 & Conceitos Básicos de Iluminação Industrial.  Lâmpadas.  Luminárias.  Dispositivos de Controle.  Iluminação de Interiores. 
\\ 
 & Divisão da Carga Elétrica em Circuitos. 
\\ 
 & Dimensionamento de Condutores. 
\\ 
 & Dimensionamento de Condutos. 
\\ 
 & Proteção Contra Sobrecorrentes. 
\\ 
 & Alimentação de Motores Elétricos. Métodos de Partida em MITs. 
\\ 
 & Sistema de Proteção Contra Descargas Atmosféricas.  
\\ 
 & Proteção Contra Choques Elétricos.  Aterramento. }

\def\PUDmetodo{Aulas expositórias que podem ser teóricas e/ou práticas, onde as práticas no laboratório serão marcadas no decorrer da disciplina. Pode também ser utilizada, dependendo do aceite ou não dos alunos, a metodologia baseada em projetos, em que é proposto um projeto de instalações elétricas no início do semestre e o seu desenvolvimento é acompanhado e avaliado periódicamente durante todo o percurso do semestre.}

\def\PUDavalia{A avaliação será feita com aplicação de provas teóricas e/ou práticas, além da possibilidade de inclusão de trabalhos, seminários e projetos de instalações elétricas no decorrer da disciplina.}

\def\PUDbibbasica{ & \href{www.biblioteca.ifce.edu.br/index.asp?codigo_sophia=23910}{Cotrim}, Ademaro A. M. B.  Instalações Elétricas.  5º ed.  São Paulo, SP: Pearson Prentice Hall, 2009.  496p.  ISBN: 9788576052081.
\\ 
 & \href{www.biblioteca.ifce.edu.br/index.asp?codigo_sophia=24178}{Creder}, Hélio.  Instalações Elétricas.  15º ed.  Rio de Janeiro, RJ: LTC, 2007.  428p.  ISBN: 9788521615675.
\\ 
 & \href{www.biblioteca.ifce.edu.br/index.asp?codigo_sophia=24294}{Mamede Filho}, João.  Instalações Elétricas Industriais.  7º ed.  Rio de Janeiro, RJ: LTC, 2007.  914p.  ISBN: 9788521615200. }

\def\PUDbibcomplementar{ & \href{www.biblioteca.ifce.edu.br/index.asp?codigo_sophia=24052}{Mamede Filho}, João.  Manual de equipamentos elétricos.  3º ed.  Rio de Janeiro, RJ: LTC, 2005.  778p.  ISBN: 9788521614364.
\\ 
 & \href{www.biblioteca.ifce.edu.br/index.asp?codigo_sophia=24525}{Prazeres}, Romildo Alves dos.  Redes de distribuição de energia elétrica e subestações.  Curitiba, PR: Base Editorial, 2010.  176p.  ISBN: 9788579055614. 
\\ 
 & \href{www.biblioteca.ifce.edu.br/index.asp?codigo_sophia=24472}{Walenia}, Paulo Sérgio.  Projetos Elétricos Industriais.  Curitiba, PR: Base Editorial, 2010.  288p.  ISBN: 9788579055577.
\\
 & \href{www.biblioteca.ifce.edu.br/index.asp?codigo_sophia=61952}{Niskier}, Júlio.  Instalações elétricas.  6º ed.  Rio de Janeiro, RJ: LTC, 2015.  443p.  ISBN: 9788521622130.
\\
 & \href{www.biblioteca.ifce.edu.br/index.asp?codigo_sophia=65036}{Filippo Filho}, Guilherme.  Motor de indução: princípios de funcionamento, características operacionais, aplicações, acionamentos e comandos.  2º ed.  São Paulo, SP: Érica, 2016.  294p.  ISBN: 9788536504483. }
%\\ 
 %&  NISKIER, Júlio.  Instalações Elétricas.  ISBN : 85-21-61589-2, 5a Edição, LTC. }

\def\PUDautor{Celso Rogério Schmidlin Júnior}

\def\PUDdata{2013-06-26}

\def\PUDrevisao{3}

\def\PUDrevisor{Celso Rogério Schmidlin Júnior}

\def\PUDdatarevisao{2019-05-15}

\PUDcorpo
